\section{Results}
\label{sec:results}

\begin{table*}[!t]
\centering
\footnotesize
\renewcommand{\arraystretch}{1.25}
\setlength{\tabcolsep}{3pt}
\begin{tabular*}{\textwidth}{@{\extracolsep{\fill}}llcccccccc@{}}
\toprule
& & \multicolumn{4}{c}{Low-Level} & \multicolumn{4}{c}{High-Level} \\
\cmidrule(lr){3-6}\cmidrule(lr){7-10}
\multicolumn{2}{l}{\textbf{Method}}
& PixCorr$\uparrow$ & SSIM$\uparrow$ & Alex(2)$\uparrow$ & Alex(5)$\uparrow$
& Incep$\uparrow$ & CLIP$\uparrow$ & Eff$\downarrow$ & SwAV$\downarrow$ \\
\midrule
\multirow{5}{*}{\shortstack{Voxel-\\based}}
& Takagi \textit{et al.}$^\ddagger$ (CVPR'23) \cite{10205187}
& 0.246 & 0.410 & 78.9\% & 85.6\%
& 83.8\% & 82.1\% & 0.811 & 0.504 \\
& Brain-Diffuser$^\ddagger$ (Sci-Rep'23) \cite{ozcelik2023brain}
& 0.273 & 0.365 & 94.4\% & 96.6\%
& 91.3\% & 90.9\% & 0.728 & 0.421 \\
& MindEye$^\ddagger$ (NeurIPS'23) \cite{scotti2023reconstructingmindseyefmritoimage}
& 0.319 & 0.360 & 92.8\% & 96.9\%
& 94.6\% & \underline{93.3\%} & 0.648 & 0.377 \\
& MindEye2$^\dagger$ (ICML'24) \cite{scotti2024mindeye2sharedsubjectmodelsenable}
& \underline{0.322} & \underline{0.431} & 96.1\% & \underline{98.6\%}
& \underline{95.4\%} & 93.0\% & \underline{0.619} & \underline{0.344} \\
& Brain-IT$^\dagger$ (arXiv'25) \cite{beliy2025brainitimagereconstructionfmri}
& \textbf{0.386} & \textbf{0.486} & \textbf{98.4\%} & \textbf{99.5\%}
& \textbf{97.3\%} & \textbf{96.4\%} & \textbf{0.564} & \textbf{0.320} \\
\midrule
\multirow{3}{*}{Flat-map}
& NeuroPictor$^\dagger$ (ECCV'24) \cite{huo2024neuropictorrefiningfmritoimagereconstruction}
& 0.229 & 0.375 & \underline{96.5\%} & 98.4\%
& 94.5\% & \underline{93.3\%} & 0.639 & 0.350 \\
& \textbf{CortexDecoder (Low-Level)}
& 0.321 & 0.393 & 94.4\% & 97.9\%
& 89.7\% & 74.3\% & 0.856 & 0.590 \\
& \textbf{CortexDecoder (Final)}
& 0.311 & 0.408 & 93.2\% & 96.1\%
& 91.4\% & 89.9\% & 0.697 & 0.404 \\
\bottomrule
\end{tabular*}
\caption{Reconstruction metrics on NSD shared1000, averaged over Subjects 1, 2, 5 and 7. Baselines are grouped by fMRI input representation and taken as reported by each paper, except $^\ddagger$, which are recomputed by MindEye2 \cite{scotti2024mindeye2sharedsubjectmodelsenable} on the full 40 sessions. \textbf{Bold}: best; \underline{underline}: second best. $^\dagger$Pretrains on multiple subjects.}
\label{tab:reconstruction}
\end{table*}


We evaluate CortexDecoder on the shared1000 from the NSD, a dataset completely separated from training and validation to prevent leakage. All metrics computed compare reconstruction against the stimulus images viewed by subject one. We report, PixCorr (Pixel Correlation), SSIM (Structural Similarity Index Measure) \cite{1284395}, AlexNet(2), AlexNet(5) \cite{NIPS2012_c399862d}, InceptionV3 \cite{szegedy2015rethinkinginceptionarchitecturecomputer}, CLIP \cite{radford2021learningtransferablevisualmodels}, EfficientNet-B1 \cite{tan2020efficientnetrethinkingmodelscaling} and SwAV-ResNet50 \cite{caron2021unsupervisedlearningvisualfeatures}. PixCorr and SSIM measure pixel and structural similarity while AlexNet, InceptionV3, CLIP measure feature space similarity. EfficientNet and SwAV report correlation distance, where lower is better.

We report metrics for both the lower-level DIP reconstruction and the final diffusion reconstruction. This allows us to evaluate the contribution of the lower-level branch separately from the semantic conditioning. Because diffusion is stochastic, we generate ten candidates per image. Then we select the generated image with the lowest pixelwise MSE to the low-level reconstruction after resizing. This approach does not use ground-truth images and favors candidates that preserve the predicted structure.

\subsection{Reconstruction Benchmarks}
CortexDecoder is competitive with prior methods in capturing low-level and mid-level features in reconstructed images. Specifically, our model performs well for PixCorr and SSIM reflecting the contribution of the lower-level branch. However, we report weaker results in terms of capturing semantic features in the reconstructed images.

From the results it is clear that the introduction of semantic conditioning negatively impacts the modeling of lower level features. For example, the low-level method on its own scores higher on the low-level benchmarks compared to the samples where semantic conditioning was introduced in the diffusion stage. We conclude that the final reconstruction which takes into account both semantic and low-level representation improves on semantic metrics at the expense of low-level performance.


